% ============================================================================
% NeuroGraph: A Temporal Knowledge Graph Operating System for AI Memory
% Comprehensive Research Paper
% ============================================================================
\documentclass[11pt,a4paper,twocolumn]{article}

% ─── Packages ───────────────────────────────────────────────────────────────
\usepackage[utf8]{inputenc}
\usepackage[T1]{fontenc}
\usepackage{lmodern}
\usepackage[margin=2cm]{geometry}
\usepackage{hyperref}
\usepackage{graphicx}
\usepackage{amsmath,amssymb}
\usepackage{algorithm}
\usepackage{algpseudocode}
\usepackage{booktabs}
\usepackage{multirow}
\usepackage{listings}
\usepackage{xcolor}
\usepackage{tikz}
\usepackage{caption}
\usepackage{subcaption}
\usepackage{enumitem}
\usepackage{fancyhdr}
\usepackage[numbers,sort&compress]{natbib}
\usepackage{float}
\usepackage{tabularx}

\usetikzlibrary{shapes,arrows,positioning,fit,backgrounds}

% ─── Code listing style ────────────────────────────────────────────────────
\definecolor{codebg}{HTML}{F5F5F5}
\definecolor{codegreen}{HTML}{2E7D32}
\definecolor{codepurple}{HTML}{7B1FA2}
\definecolor{codeorange}{HTML}{E65100}
\definecolor{codegray}{HTML}{757575}

\lstdefinestyle{rustcode}{
  backgroundcolor=\color{codebg},
  basicstyle=\ttfamily\scriptsize,
  breaklines=true,
  captionpos=b,
  commentstyle=\color{codegray}\itshape,
  keywordstyle=\color{codepurple}\bfseries,
  stringstyle=\color{codegreen},
  numberstyle=\tiny\color{codegray},
  numbers=left,
  frame=single,
  rulecolor=\color{codegray!30},
  language=C, % closest built-in to Rust
  morekeywords={fn,let,mut,pub,struct,impl,async,await,self,use,mod,
    trait,enum,match,if,else,for,in,return,Ok,Err,Some,None,
    String,Vec,Arc,HashMap,Result,Option},
  emph={NeuroGraph,IngestionPipeline,HybridRetriever,DriftSearch,
    TieredMemory,MultiGraphMemory,BranchManager,ZettelkastenMemory,
    MemoryEvolution,IntentRouter,FusionEngine,PdfParser},
  emphstyle=\color{codeorange}\bfseries,
  tabsize=4,
  showstringspaces=false,
}

\lstset{style=rustcode}

% ─── Metadata ──────────────────────────────────────────────────────────────
\title{%
  \textbf{NeuroGraph: A Temporal Knowledge Graph Operating System\\
  for AI Memory with Self-Evolving Retrieval}\\[6pt]
  \large Ingest Anything, Remember Everything, Forget Intelligently, Reason Visually
}
\author{
  Ashutosh Kumar Singh\\
  \texttt{ashutosh@neurograph.dev}\\
  NeuroGraph AI\\
  \url{https://github.com/neurographai/neurograph}
}
\date{March 2026}

% ─── Headers ───────────────────────────────────────────────────────────────
\pagestyle{fancy}
\fancyhf{}
\fancyhead[L]{\small NeuroGraph Research Paper}
\fancyhead[R]{\small March 2026}
\fancyfoot[C]{\thepage}

\begin{document}
\maketitle

% ════════════════════════════════════════════════════════════════════════════
\begin{abstract}
We present \textbf{NeuroGraph}, a production-grade knowledge graph operating
system implemented in Rust that unifies temporal knowledge management,
multi-strategy hybrid retrieval, and biologically-inspired memory evolution
within a single, ergonomic framework. NeuroGraph introduces several novel
contributions: (1)~a \emph{richly-typed graph data model} with 16-field
entities, bi-temporal relationships, episodic provenance, saga-based
conversation threading, and a dual-mode ontology system (prescribed
vs.~learned types); (2)~a \emph{seven-strategy hybrid retrieval engine}
combining semantic vector search, BM25 keyword indexing, Personalized
PageRank, graph traversal, cross-encoder reranking, Maximal Marginal
Relevance (MMR) for diversity, and node-distance graph-proximity
reranking, all fused via Reciprocal Rank Fusion (RRF); (3)~a
\emph{composable search configuration system} with per-channel
(entity/relationship/episode/community) independent method and reranker
selection, multi-channel typed results with provenance scoring, and
15+ pre-built search recipes; (4)~a \emph{two-level query engine} with
a strategy-dispatch router (Local/Global/Temporal/MultiHop) and a
MAGMA-style intent router classifying eight distinct query intents
across four orthogonal graph views (semantic, temporal, causal, entity);
(5)~a \emph{tiered memory hierarchy} (L1--L4: Working, Episodic,
Semantic, Procedural) with automatic promotion, demotion, and
Ebbinghaus-curve forgetting; (6)~an \emph{RL-guided memory evolution
engine} using Q-table reinforcement learning for keep/forget decisions;
(7)~a \emph{graph-level forgetting engine} with exponential importance
decay, access-frequency boosting, and TTL-based expiration; (8)~a
\emph{Git-like branch isolation protocol} with copy-on-write deltas,
diffing, and merging; (9)~a \emph{provider-agnostic LLM abstraction}
with a dyn-compatible trait, OpenAI-format universal client, SHA-256
keyed response caching with LRU eviction, and per-prompt-type token
usage tracking; and (10)~a \emph{research paper intelligence platform}
with PDF ingestion, multi-source academic search, RAG chat, and Model
Context Protocol (MCP) integration for Claude, Cursor, and VS~Code
Copilot. The system is architected as a Rust workspace of four crates
totaling over 35,000 lines of code, with zero-configuration defaults,
pluggable storage backends (in-memory DashMap, embedded Sled, planned
Neo4j), and a REST API with an embedded visualization dashboard. We
describe the architecture, algorithms, and implementation in detail.
\end{abstract}

\textbf{Keywords:} Knowledge Graph, Temporal Reasoning, Graph RAG,
Hybrid Retrieval, Memory Architecture, Reinforcement Learning,
Multi-Agent Systems, Rust

% ════════════════════════════════════════════════════════════════════════════
\section{Introduction}
\label{sec:intro}

The proliferation of large language models (LLMs) has exposed a
fundamental limitation: LLMs lack persistent, structured memory.
While retrieval-augmented generation (RAG) systems partially address
this by grounding responses in external documents, they treat knowledge
as flat text chunks, discarding the relational structure that makes
knowledge \emph{actionable}.

Knowledge graphs (KGs) offer a solution by representing information as
typed entities connected by labeled relationships. However, existing KG
systems suffer from several deficiencies: (1)~they lack temporal
awareness---facts are treated as eternally valid; (2)~they employ
single-strategy retrieval that fails on diverse query types;
(3)~they have no mechanism for intelligent forgetting, leading to
unbounded memory growth; and (4)~they provide no isolation mechanism
for concurrent multi-agent access.

NeuroGraph addresses all four deficiencies through an integrated
operating-system metaphor: knowledge is \emph{ingested} through a
multi-stage pipeline, \emph{stored} in a temporal graph with
bi-temporal validity intervals, \emph{retrieved} through a
five-strategy hybrid engine, \emph{evolved} through RL-guided
forgetting, and \emph{isolated} through Git-like branching.

\subsection{Design Philosophy}

NeuroGraph follows three core principles:

\begin{enumerate}[leftmargin=*]
  \item \textbf{Zero-Configuration Start}: A single line---
    \lstinline{NeuroGraph::builder().build().await?}---creates a
    fully functional instance with in-memory storage, local hash
    embeddings, and regex-based entity extraction. No API keys,
    no external services.
  \item \textbf{Progressive Enhancement}: Each subsystem (LLM
    extraction, tiered memory, multi-graph memory, evolution) is
    independently opt-in via the builder pattern.
  \item \textbf{Rust-Native Performance}: The entire system is
    implemented in safe Rust with async/await concurrency via
    Tokio, lock-free data structures (DashMap, parking\_lot),
    and zero-copy serialization.
\end{enumerate}

\subsection{Contributions}

This paper makes the following contributions:

\begin{itemize}[leftmargin=*]
  \item A \textbf{complete system architecture} for a temporal
    knowledge graph OS with pluggable drivers, embedders, and LLM
    clients (\S\ref{sec:architecture}).
  \item A \textbf{richly-typed graph data model} with bi-temporal
    entities, provenance episodes, dual-mode ontology, and schema
    registry (\S\ref{sec:datamodel}).
  \item A \textbf{pluggable driver trait} with 25 async methods and
    two built-in implementations: lock-free in-memory (DashMap) and
    embedded persistent (Sled B-tree) (\S\ref{sec:drivers}).
  \item A \textbf{provider-agnostic LLM abstraction} with object-safe
    trait, OpenAI-format universal client, and per-call cost tracking
    (\S\ref{sec:llm}).
  \item A \textbf{two-level query engine} with strategy dispatch,
    token-budgeted context assembly, and per-query cost enforcement
    (\S\ref{sec:engine}).
  \item A \textbf{multi-stage ingestion pipeline} with extraction,
    ontology validation, embedding-based deduplication, and temporal
    conflict resolution (\S\ref{sec:ingestion}).
  \item A \textbf{seven-strategy hybrid retrieval engine} with RRF fusion,
    MMR diversity reranking, node-distance graph-proximity reranking,
    episode-mentions citation reranking, and DRIFT adaptive local/global
    search (\S\ref{sec:retrieval}).
  \item A \textbf{composable search configuration system} with per-channel
    method and reranker selection, multi-channel typed results, and 15+
    pre-built search recipes (\S\ref{sec:searchconfig}).
  \item A \textbf{biologically-inspired memory system} combining tiered
    hierarchy, Zettelkasten self-organization, and RL-guided evolution
    (\S\ref{sec:memory}).
  \item A \textbf{multi-graph memory architecture} with intent-aware
    routing across four orthogonal graph views (\S\ref{sec:multigraph}).
  \item A \textbf{temporal subsystem} with bi-temporal validity,
    point-in-time snapshots, Git-like branching, hybrid logical clock,
    fact version chains, and graph-level forgetting
    (\S\ref{sec:temporal}).
  \item A \textbf{research paper intelligence platform} with PDF
    parsing, multi-source search, RAG chat, and MCP integration
    (\S\ref{sec:research}, \S\ref{sec:mcp}).
\end{itemize}

% ════════════════════════════════════════════════════════════════════════════
\section{System Architecture}
\label{sec:architecture}

\subsection{Workspace Organization}

NeuroGraph is organized as a Rust Cargo workspace comprising four
crates:

\begin{table}[H]
\centering
\caption{NeuroGraph workspace crate organization.}
\label{tab:crates}
\begin{tabularx}{\columnwidth}{lX}
\toprule
\textbf{Crate} & \textbf{Responsibility} \\
\midrule
\texttt{neurograph-core} & Core library: graph, ingestion, retrieval,
  memory, temporal, PDF, papers, chat, server, embedders, community \\
\texttt{neurograph-cli} & Command-line interface with 12 subcommands \\
\texttt{neurograph-mcp} & Model Context Protocol server (stdio/SSE) \\
\texttt{neurograph-eval} & Evaluation harness and benchmarks \\
\bottomrule
\end{tabularx}
\end{table}

The core crate exposes 17 top-level modules, each encapsulating a
distinct subsystem. The workspace uses Rust 2021 edition with a
minimum supported Rust version (MSRV) of 1.82.

\subsection{Layered Architecture}

The system follows a strict layered architecture:

\begin{enumerate}[leftmargin=*]
  \item \textbf{Public API Layer} (\texttt{NeuroGraph} struct): Exposes
    \texttt{add\_text()}, \texttt{add\_json()}, \texttt{query()},
    \texttt{at()}, \texttt{entity\_history()}, \texttt{what\_changed()},
    \texttt{remember()}, \texttt{recall()}, \texttt{detect\_communities()},
    and builder methods.
  \item \textbf{Engine Layer}: Query routing
    (\texttt{QueryRouter}), ingestion pipeline
    (\texttt{IngestionPipeline}), community detection.
  \item \textbf{Service Layer}: LLM clients (OpenAI, Ollama),
    embedding providers (OpenAI, Ollama, FastEmbed/Hash, cached),
    paper search APIs.
  \item \textbf{Driver Layer}: Pluggable storage backends via the
    \texttt{GraphDriver} trait---\texttt{MemoryDriver} (in-memory
    with \texttt{DashMap}), \texttt{EmbeddedDriver} (Sled persistent
    B-tree), and planned Neo4j driver.
\end{enumerate}

\subsection{Builder Pattern and Configuration}

The \texttt{NeuroGraphBuilder} provides a fluent API for progressive
configuration:

\begin{lstlisting}
let ng = NeuroGraph::builder()
    .name("research-kg")
    .embedded("./data")
    .budget(5.0)            // USD limit
    .multigraph()           // Enable MAGMA
    .tiered_memory()        // Enable L1-L4
    .evolution()            // Enable RL forgetting
    .build().await?;
\end{lstlisting}

The \texttt{.full()} shorthand enables all v0.2 subsystems
simultaneously. When no configuration is provided,
\texttt{NeuroGraphConfig::zero\_config()} creates sensible defaults:
in-memory storage, local hash embeddings (128-dimensional), no LLM
(regex extraction), and a \$10 USD budget ceiling.

\subsection{Error Model and Cost Tracking}

NeuroGraph employs a typed error hierarchy via \texttt{thiserror}:

\begin{itemize}[leftmargin=*]
  \item \texttt{NeuroGraphError::Driver} --- storage failures
  \item \texttt{NeuroGraphError::Llm} --- LLM API errors
  \item \texttt{NeuroGraphError::Embedder} --- embedding failures
  \item \texttt{NeuroGraphError::BudgetExceeded\{spent, limit\}} ---
    hard budget enforcement
  \item \texttt{NeuroGraphError::Temporal} --- date parsing, snapshot
  \item \texttt{NeuroGraphError::Community} --- detection/summarization
\end{itemize}

Every LLM call records its cost via the \texttt{CostTracker}, and
the \texttt{ConcurrencyLimiter} throttles parallel LLM calls to
prevent rate-limit breaches.

\subsection{Graph Data Model}
\label{sec:datamodel}

The \texttt{graph} module defines the core data structures that
underpin all subsystems. Each type uses UUID~v4 identifiers,
\texttt{serde} serialization, and the builder pattern.

\paragraph{Entity.}
The \texttt{Entity} struct carries 16 fields: \texttt{id} (UUID),
\texttt{name}, \texttt{entity\_type} (wrapped \texttt{EntityType(String)}
for type safety), \texttt{summary} (LLM-generated regional
description), \texttt{name\_embedding} (\texttt{Option<Vec<f32>>}),
\texttt{group\_id} (multi-tenant partition), \texttt{labels}
(flexible categorization, e.g., \texttt{["Entity", "Person"]}),
\texttt{attributes} (\texttt{HashMap<String, serde\_json::Value>}),
\texttt{importance\_score} ($[0,1]$, PageRank + access frequency),
\texttt{access\_count}, \texttt{community\_ids}
(\texttt{HashMap<u32, String>} mapping hierarchy level to community),
\texttt{created\_at}, \texttt{updated\_at}, and \texttt{metadata}.
The \texttt{record\_access()} method atomically increments the
access counter and touches the modification timestamp.

\paragraph{Relationship.}
Edges carry bi-temporal fields directly from Graphiti's
\texttt{EntityEdge}: \texttt{valid\_from} / \texttt{valid\_until}
(real-world validity), \texttt{created\_at} / \texttt{expired\_at}
(system time). Additional fields include \texttt{fact} (natural
language statement), \texttt{fact\_embedding}, \texttt{weight}
(default 1.0), \texttt{confidence} ($[0,1]$), and
\texttt{episode\_ids} (provenance chain). The \texttt{is\_valid\_at(t)}
method enables point-in-time queries, and \texttt{invalidate()}
sets both \texttt{valid\_until} and \texttt{expired\_at}.

\paragraph{Episode.}
Episodes are immutable provenance records with an \texttt{EpisodeType}
enum supporting six source types: \texttt{Message}, \texttt{Json},
\texttt{Text}, \texttt{Image}, \texttt{Audio}, and \texttt{Pdf}.
Each episode carries \texttt{valid\_at} (real-world timestamp of
the source data) and \texttt{entity\_edge\_ids} (extracted
relationship UUIDs).

\paragraph{Saga (Conversation Thread).}
The \texttt{Saga} struct groups \texttt{Episode}s into ordered
conversation threads, closing a critical competitive gap with
Graphiti's \texttt{SagaNode}. Each saga maintains: \texttt{SagaId}
(UUID), \texttt{name}, optional \texttt{description} and
\texttt{source} identifiers, \texttt{group\_id} for multi-tenant
isolation, and an append-only \texttt{episode\_ids} vector that
preserves chronological insertion order.

The graph structure uses two well-known edge types:
\texttt{HAS\_EPISODE} (with a \texttt{position} attribute) from
Saga to Episode, and \texttt{NEXT\_EPISODE} between consecutive
episodes. The \texttt{episode\_pairs()} method generates
$(\text{source}, \text{target})$ pairs for constructing the
linked-list threading edges. This enables queries such as ``What
was discussed in conversation~X?'' and chronological replay of
multi-turn interactions.

\paragraph{TemporalFact.}
The \texttt{TemporalFact} struct wraps a relationship with a
dedicated \texttt{TemporalValidity} struct (\texttt{valid\_from},
\texttt{valid\_until}, \texttt{recorded\_at},
\texttt{invalidated\_at}), plus \texttt{superseded\_by}
(\texttt{Option<RelationshipId>}) and
\texttt{invalidation\_reason}. This enables fact version chains
where superseded facts link to their replacements.

\paragraph{Ontology.}
The \texttt{Ontology} system supports two modes:
\emph{prescribed} (user-defined \texttt{EntityTypeDefinition}
with expected attributes, data types, required flags, and
few-shot examples) and \emph{learned} (auto-discovered types
from extraction). \texttt{RelationshipTypeDefinition} specifies
allowed source/target entity types. The \texttt{to\_extraction\_prompt()}
method generates LLM-ready markdown from the ontology, and
\texttt{strict\_mode} constrains extraction to prescribed types only.

\paragraph{GraphSchema.}
The \texttt{GraphSchema} registry tracks all types currently present
in the graph: entity type counts, relationship type counts, known
type sets, and totals. It auto-registers learned types into the
ontology via \texttt{record\_entity\_type()} and
\texttt{record\_relationship\_type()}.

\subsection{GraphDriver Trait and Implementations}
\label{sec:drivers}

The \texttt{GraphDriver} trait defines 25 \texttt{async} methods
across six operation categories:

\begin{table}[H]
\centering
\caption{GraphDriver trait method categories.}
\label{tab:driver}
\begin{tabularx}{\columnwidth}{lcX}
\toprule
\textbf{Category} & \textbf{Methods} & \textbf{Operations} \\
\midrule
Entity & 5 & store, get, batch get, delete, list \\
Relationship & 4 & store, batch store, get, get by entity \\
Episode & 3 & store, get, list recent \\
Community & 5 & store, get, by level, dirty, list \\
Search & 3 & vector entities, vector rels, text search \\
Graph/Temporal & 3 & traverse, snapshot\_at, stats \\
\bottomrule
\end{tabularx}
\end{table}

Two implementations are provided:

\begin{itemize}[leftmargin=*]
  \item \textbf{MemoryDriver}: Uses \texttt{DashMap<K, V>} for
    lock-free concurrent access to entities, relationships,
    episodes, and communities. Vector search uses brute-force
    cosine similarity ($O(nk)$); text search uses case-insensitive
    substring matching. Temporal filtering applies
    \texttt{is\_valid\_at()} to all relationships in
    \texttt{snapshot\_at()}.
  \item \textbf{EmbeddedDriver}: Uses Sled embedded B-tree with
    separate trees per data type (\texttt{entities},
    \texttt{relationships}, \texttt{episodes}, \texttt{communities}).
    Serialization via \texttt{bincode} for compact binary storage.
    Inherits the same vector/text search algorithms as MemoryDriver
    but reads from persistent disk.
\end{itemize}

\texttt{GraphDriverConfig} includes \texttt{storage\_path},
\texttt{max\_connections}, and optional \texttt{uri}/\texttt{username}/
\texttt{password} fields for planned remote drivers (Neo4j, FalkorDB).

\subsection{LLM Abstraction Layer}
\label{sec:llm}

The \texttt{llm} module provides a provider-agnostic LLM interface
with response caching and per-prompt-type token tracking.

\paragraph{LlmClient Trait.}
The core trait defines three methods: \texttt{model\_name()},
\texttt{cost\_per\_token()} (returns \texttt{(input\_cost, output\_cost)}
in USD), and \texttt{complete(request) -> CompletionResponse}.
The trait is \texttt{dyn}-compatible (object-safe); structured
JSON deserialization is provided via a free function
\texttt{complete\_structured<T>()} rather than a generic trait method,
preserving dynamic dispatch.

\paragraph{CompletionRequest.}
Encapsulates \texttt{system\_prompt}, \texttt{user\_prompt},
\texttt{temperature}, \texttt{max\_tokens}, and \texttt{json\_mode}
(requests \texttt{response\_format: json\_object}).

\paragraph{GenericLlmClient.}
Uses raw \texttt{reqwest} HTTP calls to any OpenAI-compatible
endpoint (Ollama at \texttt{localhost:11434/v1}, vLLM, LiteLLM,
Azure OpenAI). Features include: configurable timeout, bearer
token auth, rate-limit detection (HTTP 429 $\to$
\texttt{LlmError::RateLimited}), and per-call cost calculation
using \texttt{LlmConfig::calculate\_cost(input\_tokens,
output\_tokens)}.

\paragraph{LlmConfig Presets.}
Pre-built configurations: \texttt{gpt4o\_mini()} ($\$0.15/\$0.60$
per million tokens), \texttt{ollama(model)} (zero-cost, local
endpoint). Each config stores \texttt{input\_cost\_per\_million}
and \texttt{output\_cost\_per\_million} for budget tracking.

\paragraph{Response Cache (\texttt{LlmCache}).}
Prevents redundant API calls for identical extraction prompts---closing
a gap where GraphRAG provides a modular caching factory. Cache keys are
computed deterministically via SHA-256 over the concatenation of prompt
and model name. The cache supports:
\begin{itemize}[leftmargin=*]
  \item \textbf{LRU Eviction}: When the cache exceeds
    \texttt{max\_entries} (configurable, default 1,000), the oldest
    entry by \texttt{cached\_at} timestamp is evicted.
  \item \textbf{TTL Expiration}: Each \texttt{CachedResponse} carries
    an optional TTL (\texttt{chrono::Duration}). Expired entries are
    treated as cache misses and evicted on the next write pass.
  \item \textbf{Observability}: \texttt{CacheStats} tracks hits,
    misses, and evictions, with a \texttt{hit\_rate()} convenience
    method. Thread-safety is provided via \texttt{tokio::sync::RwLock}.
\end{itemize}
Each cached entry records the response content, model name, input/output
token counts, and timestamp---enabling cost tracking even for cached
results.

\paragraph{Token Usage Tracker (\texttt{TokenTracker}).}
Tracks LLM token consumption at granular, per-prompt-type level.
Eight built-in \texttt{PromptType} variants categorize API calls:
\texttt{EntityExtraction}, \texttt{RelationshipExtraction},
\texttt{EntityResolution}, \texttt{CommunitySummary},
\texttt{QueryRewrite}, \texttt{AnswerGeneration}, \texttt{Reranking},
and \texttt{ConflictResolution}, plus \texttt{Custom(String)} for
extension. Each \texttt{record()} call updates both per-type
counters (\texttt{TokenUsage}: input/output tokens, call count,
estimated cost) and global atomic counters
(\texttt{AtomicU64} for lock-free reads). The \texttt{print\_summary()}
method renders a formatted table for cost analysis, and
\texttt{get\_summary()} returns the full
\texttt{HashMap<PromptType, TokenUsage>} for programmatic access.

\subsection{Query Engine}
\label{sec:engine}

The \texttt{engine} module orchestrates query execution through
a two-level routing architecture:

\paragraph{Level 1: QueryRouter (Strategy Dispatch).}
Classifies queries into four types via keyword heuristic scoring:
\texttt{Local} (entity-focused, default), \texttt{Global}
(dataset-wide: ``summarize'', ``overview'', ``all''),
\texttt{Temporal} (``when'', ``before'', ``history'', ``in 20XX''),
and \texttt{MultiHop} (``connected'', ``related'', ``path'').
This router dispatches to the appropriate \texttt{QueryStrategy}
implementation:

\begin{itemize}[leftmargin=*]
  \item \textbf{LocalStrategy}: Hybrid retrieval $\to$ relationship
    fetch $\to$ context assembly $\to$ LLM answer generation.
    Returns raw context if no LLM is configured.
  \item \textbf{GlobalStrategy}: Community map-reduce pipeline---
    (1) retrieve community summaries, (2) map: LLM partial answer
    per community, (3) reduce: synthesize final answer. Falls back
    to hybrid search when no communities exist.
\end{itemize}

\paragraph{Level 2: IntentRouter (Graph View Selection).}
When multi-graph memory is enabled, the \texttt{IntentRouter}
(\S\ref{sec:multigraph}) performs finer-grained 8-type intent
classification and selects which graph views to query. The two
routers are complementary: the engine router selects the execution
strategy, while the intent router selects the data sources.

\paragraph{ContextBuilder.}
Assembles LLM prompts from graph data within a configurable token
budget (default 4,000 tokens, estimated at $\approx$4
chars/token). Content is added in priority order: (1) entity
descriptions with types and summaries, (2) relationship facts
with temporal validity markers (\texttt{\checkmark~current} vs.
\texttt{\texttimes~historical}), (3) community summaries for
global context. Each section is truncated at the budget boundary.

\paragraph{QueryBudget.}
Per-query cost enforcement using atomic microdollar tracking
(\texttt{AtomicU64} storing hundredths of cents). Provides
\texttt{record\_cost()}, \texttt{can\_afford()}, and \texttt{reset()}
operations.

\subsection{Utility Infrastructure}
\label{sec:utils}

\paragraph{ConcurrencyLimiter.}
Semaphore-based rate limiter (from Graphiti's
\texttt{semaphore\_gather()} pattern) wrapping \texttt{tokio::sync::%
Semaphore} with an atomic active-count tracker. Default limit: 10
concurrent LLM calls.

\paragraph{CostTracker.}
Process-wide atomic cost accumulator storing microdollars
(\texttt{Arc<AtomicU64>}). Provides \texttt{record()},
\texttt{total\_cost\_usd()}, \texttt{is\_over\_budget()}, and
\texttt{remaining\_usd()} for budget enforcement across
all operations.

\paragraph{Hashing.}
Two-phase deduplication hashing from Graphiti:
(1) deterministic SHA-256 content hash with text normalization
(lowercase, whitespace collapse); (2) BLAKE3 fast hash for
high-throughput scenarios. Helper functions
\texttt{entity\_dedup\_key(name, type)} and
\texttt{relationship\_dedup\_key(source, target, type)} generate
consistent dedup keys.

\paragraph{Text Utilities.}
\texttt{normalize\_name()} (lowercase, replace underscores/hyphens),
\texttt{extract\_terms()} (stopword-filtered tokenization),
\texttt{truncate\_to\_tokens()} (word-boundary-aware truncation),
\texttt{estimate\_tokens()} ($\lceil\text{chars}/4\rceil$),
and \texttt{chunk\_text()} (sentence-boundary-respecting chunking).

% ════════════════════════════════════════════════════════════════════════════
\section{Ingestion Pipeline}
\label{sec:ingestion}

\subsection{Pipeline Architecture}

The ingestion pipeline transforms raw input (text, JSON, PDF) into
a structured knowledge graph through six stages:

\begin{equation}
\text{Src} \xrightarrow{1} \text{Ep}
  \xrightarrow{2} \text{Ext}
  \xrightarrow{3} \text{Val}
  \xrightarrow{4} \text{Dup}
  \xrightarrow{5} \text{Cfl}
  \xrightarrow{6} \text{Str}
\end{equation}

\paragraph{Stage 1: Episode Creation.}
Every ingestion creates an \texttt{Episode}---a provenance record
linking the raw input to the entities and relationships it produces.
Episodes carry a \texttt{group\_id} for multi-tenant isolation and
a \texttt{created\_at} timestamp for temporal ordering.

\paragraph{Stage 2: Entity/Relationship Extraction.}
The \texttt{TextExtractor} supports two modes: (a)~LLM-based
extraction using structured prompts to an OpenAI-compatible API,
returning typed entities with summaries and relationships with
confidence scores; (b)~regex-only fallback that identifies
capitalized proper nouns and verb-mediated relationships, requiring
no API keys.

\paragraph{Stage 3: Ontology Validation.}
The \texttt{SchemaValidator} checks extracted entities and relationships
against a configurable \texttt{Ontology} that defines allowed entity
types, relationship types, and constraints. Invalid items are logged
and skipped; warnings are emitted for items requiring normalization.

\paragraph{Stage 4: Embedding-Based Deduplication.}
The \texttt{Deduplicator} prevents entity proliferation by computing
cosine similarity between a new entity's name embedding and all
existing entities of the same type. If similarity exceeds a threshold
(default 0.85), the entities are merged:

\begin{equation}
\text{sim}(\mathbf{e}_\text{new}, \mathbf{e}_\text{existing}) > \tau
\implies \text{Merge}(\mathbf{e}_\text{existing}, \mathbf{e}_\text{new})
\end{equation}

Merge combines summaries and preserves the existing entity's ID,
ensuring referential integrity.

\paragraph{Stage 5: Temporal Conflict Resolution.}
The \texttt{ConflictResolver} detects four conflict types between
a new relationship and existing ones sharing the same source entity
and relationship type:

\begin{table}[H]
\centering
\caption{Conflict types and resolution strategies.}
\label{tab:conflicts}
\begin{tabularx}{\columnwidth}{lX}
\toprule
\textbf{Type} & \textbf{Resolution} \\
\midrule
Redundant & Skip (identical fact already exists) \\
Contradiction & Invalidate old by setting
  \texttt{valid\_until = now} \\
Supersession & Same as contradiction (newer fact wins) \\
SoftConflict & Coexist (non-exclusive relationship) \\
\bottomrule
\end{tabularx}
\end{table}

\paragraph{Stage 6: Embedding and Storage.}
Each surviving entity receives a name embedding and each relationship
a fact embedding. Both are stored via the \texttt{GraphDriver} trait.

\subsection{JSON Ingestion}

The \texttt{JsonExtractor} recursively walks JSON objects, creating
entities from keys and relationships from nesting structure. The
extracted graph is then routed through the same dedup/conflict/store
pipeline as text extraction.

% ════════════════════════════════════════════════════════════════════════════
\section{Hybrid Retrieval Engine}
\label{sec:retrieval}

\subsection{Seven-Strategy Pipeline}

NeuroGraph's retrieval engine combines seven complementary strategies
in a single pipeline, each addressing different aspects of relevance:

\begin{enumerate}[leftmargin=*]
  \item \textbf{Semantic Search}: Embeds the query via the configured
    embedder, then performs cosine similarity search over entity name
    embeddings using the driver's vector index. Default weight: 0.40.
  \item \textbf{BM25 Keyword Search}: A full inverted index with
    TF-IDF scoring, supporting configurable $k_1$ (1.2) and $b$ (0.75)
    parameters. Default weight: 0.25.
  \item \textbf{Personalized PageRank (PPR)}: Seeded from the top-5
    semantic results, PPR builds a local adjacency subgraph from
    1-hop neighborhoods and runs power iteration (damping $\alpha=0.85$,
    max iterations 20, convergence $\epsilon=10^{-8}$). Default
    weight: 0.20.
  \item \textbf{Graph Traversal}: BFS from user-specified seed
    entities up to 2 hops, scoring by hop distance with exponential
    decay. Default weight: 0.15.
  \item \textbf{Cross-Encoder Reranking}: Either rule-based (zero-cost,
    offline) or API-based, reranking the fused candidate list for
    precision. Rule-based scoring combines query-document overlap,
    entity type boosting, summary length normalization, and original
    score blending.
  \item \textbf{MMR Reranking}: Maximal Marginal Relevance
    (\S\ref{sec:mmr}) for diversity-aware result selection.
  \item \textbf{Node Distance Reranking}: Graph proximity scoring
    (\S\ref{sec:nodedist}) via BFS-based shortest-path distances.
\end{enumerate}

\subsection{Reciprocal Rank Fusion}

Results from the retrieval stages are merged using Reciprocal Rank
Fusion (RRF):

\begin{equation}
\text{RRF}(d) = \sum_{s \in \mathcal{S}} \frac{w_s}{k + r_s(d) + 1}
\label{eq:rrf}
\end{equation}

where $\mathcal{S}$ is the set of strategies, $w_s$ is the strategy
weight (summing to 1.0), $k$ is the RRF constant (default 60), and
$r_s(d)$ is the rank of document $d$ in strategy $s$'s result list.
Entities appearing in multiple lists receive additive score boosts,
naturally promoting cross-strategy consensus.

\subsection{Maximal Marginal Relevance (MMR)}
\label{sec:mmr}

MMR reranking balances relevance to the query with diversity among
selected results---closing a critical competitive gap, as both Graphiti
and GraphRAG support MMR-based reranking. The \texttt{MmrReranker}
greedily selects items using:

\begin{equation}
\text{MMR}(d) = \lambda \cdot \text{sim}(d, q)
  - (1-\lambda) \cdot \max_{d_j \in S} \text{sim}(d, d_j)
\label{eq:mmr}
\end{equation}

where $\lambda \in [0, 1]$ controls the relevance-diversity trade-off:
$\lambda = 1.0$ yields pure relevance (identical to top-$k$ by score),
$\lambda = 0.0$ maximizes diversity, and $\lambda = 0.7$ is the
recommended default. The similarity metric is configurable between
cosine similarity (default) and raw dot product.

The algorithm precomputes query-candidate similarities, then performs
$k$ greedy selection steps, each computing the maximum similarity
between each unselected candidate and all already-selected items.
Complexity is $O(k \cdot n \cdot d)$ where $n$ is the candidate
count and $d$ is the embedding dimension.

\subsection{Node Distance Reranker}
\label{sec:nodedist}

The \texttt{NodeDistanceReranker} reranks results by graph distance
from a center (focal) node, closing a gap where Graphiti has
\texttt{NodeReranker.node\_distance} and
\texttt{EdgeReranker.node\_distance}. Scores follow exponential
decay:

\begin{equation}
\text{score}(v) = \gamma^{d(v, c)}
\end{equation}

where $\gamma$ is the decay factor (default 0.5) and $d(v,c)$ is
the shortest BFS distance from node $v$ to center $c$, capped at
\texttt{max\_depth} (default 3). A \texttt{combine\_scores()} method
blends distance scores with existing relevance scores:
$\alpha \cdot s_\text{relevance} + (1-\alpha) \cdot s_\text{distance}$.

\subsection{Episode Mentions Reranker}
\label{sec:epmention}

The \texttt{EpisodeMentionsReranker} reranks results by how many
episodes (provenance sources) reference them---a proxy for
importance by citation frequency. Items mentioned across many
episodes are likely more significant. Scores are normalized:
the most-mentioned item receives 1.0, with others proportional.
A \texttt{combine\_scores()} variant blends mention counts with
relevance via weight $\alpha$. Any type implementing the
\texttt{HasEpisodeIds} trait (providing \texttt{episode\_ids()}
and \texttt{episode\_count()}) can be reranked.

\subsection{DRIFT: Dynamic Reasoning and Inference with Flexible Traversal}

DRIFT adaptively switches between local and global search strategies
based on query characteristics. Query breadth is estimated via a
composite heuristic:

\begin{multline}
\beta(q) = 0.7 \cdot \frac{|\text{broad\_kw}(q)|}{|\text{broad\_kw}(q)|
  + |\text{specific\_kw}(q)|} \\
  + \; 0.3 \cdot \frac{|q|_{\text{words}}}{15}
\end{multline}

\begin{itemize}[leftmargin=*]
  \item $\beta > \tau$: \textbf{Global} --- rank community summaries
    by keyword overlap, collect entities from top-$k$ communities,
    generate follow-up queries from key entities, refine locally.
  \item $\beta < 1 - \tau$: \textbf{Local} --- BFS from query-matching
    seed entities with per-hop score decay of 0.7.
  \item Otherwise: \textbf{Adaptive} --- execute both strategies,
    merge with dynamic weighting (local 0.6, global 0.4 when local
    finds results).
\end{itemize}

\subsection{Composable Search Configuration}
\label{sec:searchconfig}

The \texttt{SearchConfig} system enables per-channel independent
configuration of search methods and rerankers, modeled after
Graphiti's composable \texttt{SearchConfig} but made Rust-idiomatic
with the builder pattern.

\paragraph{Per-Channel Configuration.}
Each of the four channels---entities, relationships, episodes,
communities---has an independent \texttt{ChannelConfig} specifying:
\begin{itemize}[leftmargin=*]
  \item \texttt{enabled}: whether the channel participates
  \item \texttt{methods}: one or more \texttt{SearchMethod}s
    (Cosine, Bm25, Bfs, Ppr, DriftAdaptive)
  \item \texttt{reranker}: the \texttt{RerankerType}
    (Rrf, Mmr\{$\lambda$\}, CrossEncoder\{model\},
    NodeDistance\{center\_id\}, EpisodeMentions, None)
  \item \texttt{limit}: per-channel result cap
  \item \texttt{sim\_min\_score}: minimum similarity threshold
\end{itemize}

\paragraph{Builder Pattern.}
The \texttt{SearchConfigBuilder} allows selective overrides:
\begin{lstlisting}
let config = SearchConfig::builder()
    .entities(ChannelConfig {
        methods: vec![Cosine, Bm25],
        reranker: RerankerType::Mmr { lambda: 0.7 },
        limit: 20,
        ..Default::default()
    })
    .communities(ChannelConfig::disabled())
    .build();
\end{lstlisting}

\paragraph{Pre-Built Recipes.}
The \texttt{SearchRecipes} module provides 15+ pre-configured
\texttt{SearchConfig}s organized by use case
(Table~\ref{tab:recipes}).

\begin{table}[H]
\centering
\scriptsize
\caption{Pre-built search configuration recipes.}
\label{tab:recipes}
\begin{tabularx}{\columnwidth}{lX}
\toprule
\textbf{Recipe} & \textbf{Description} \\
\midrule
\texttt{combined\_hybrid\_rrf} & All channels, RRF \\
\texttt{combined\_hybrid\_mmr} & All channels, MMR diversity \\
\texttt{combined\_hybrid\_cross\_enc} & All channels, LLM reranking \\
\texttt{entity\_hybrid\_rrf} & Entity-only, RRF \\
\texttt{entity\_hybrid\_mmr} & Entity-only, MMR \\
\texttt{entity\_hybrid\_ep\_mentions} & Entity-only, citation rerank \\
\texttt{entity\_only\_cosine} & Cosine only, no reranking \\
\texttt{rel\_hybrid\_rrf} & Relationship-only, RRF \\
\texttt{rel\_hybrid\_node\_distance} & Rel-only, graph proximity \\
\texttt{graph\_aware\_ppr} & PPR+BFS (unique) \\
\texttt{drift\_adaptive} & DRIFT+MMR (unique) \\
\texttt{neighborhood\_search} & BFS, distance reranking \\
\texttt{conversation\_replay} & Episode, BM25+cosine \\
\texttt{community\_hybrid\_rrf} & Community-only, RRF \\
\texttt{community\_hybrid\_mmr} & Community-only, MMR \\
\bottomrule
\end{tabularx}
\end{table}

\subsection{Multi-Channel Search Results}
\label{sec:searchresults}

The \texttt{SearchResults} struct provides typed, multi-channel
result output---closing a gap where Graphiti returns separate
lists for edges, nodes, episodes, and communities while NeuroGraph
previously returned a flat entity list.

Each channel carries a \texttt{Vec<ScoredItem<T>>} where
\texttt{ScoredItem} attaches the original search score, source
method name, and optional reranker score. The
\texttt{effective\_score()} method returns the reranker score when
available, falling back to the original. \texttt{SearchMetadata}
records the original query, config name, candidate count, search
duration, and channels searched.

For backward compatibility, \texttt{SearchResults::flatten()}
interleaves all channels into a single \texttt{Vec<FlatResult>}
sorted by effective score descending, where each \texttt{FlatResult}
carries an id, name, \texttt{Channel} enum variant, score, and
snippet.

% ════════════════════════════════════════════════════════════════════════════
\section{Memory Architecture}
\label{sec:memory}

NeuroGraph implements a three-layer memory architecture inspired by
cognitive science: tiered storage, self-organizing Zettelkasten, and
RL-guided evolution.

\subsection{Tiered Memory (L1--L4)}

The tiered memory system maintains four levels, each with distinct
characteristics:

\begin{table}[H]
\centering
\caption{Tiered memory hierarchy.}
\label{tab:tiers}
\begin{tabularx}{\columnwidth}{clcX}
\toprule
\textbf{Tier} & \textbf{Name} & \textbf{Cap} & \textbf{Role} \\
\midrule
L1 & Working & 100 & Current context window, hot cache \\
L2 & Episodic & 10K & Recent interactions, time-indexed \\
L3 & Semantic & $\infty$ & Knowledge graph facts \\
L4 & Procedural & $\infty$ & Learned rules, heuristics \\
\bottomrule
\end{tabularx}
\end{table}

\paragraph{Promotion Policy.} Items are promoted based on access
frequency and importance score:
\begin{itemize}[leftmargin=*]
  \item L1 $\to$ L2: after 3 accesses
  \item L2 $\to$ L3: after 10 accesses \emph{and} importance $\geq$ 0.6
  \item L3 $\to$ L4: after 50 accesses (confirmed pattern)
\end{itemize}

\paragraph{Demotion and Eviction.} Items not accessed within a
configurable window (default 168 hours / 1 week) with importance
$<$ 0.3 are evicted from L2. L1 uses FIFO eviction when at capacity,
with evictees demoted to L2.

\paragraph{Search Tier Boosting.} When searching across tiers,
matches in higher tiers receive multiplicative boosts: L1 $\times$4.0,
L2 $\times$2.0, L3 $\times$1.5, L4 $\times$1.0, ensuring that
recently active memories surface first.

\paragraph{Procedural Rules.} L4 stores \texttt{LearnedRule} objects
with a pattern string, confidence score, usage counter, and source
memory IDs. During query context assembly, applicable rules are
retrieved via keyword matching and included in the LLM prompt.

\subsection{Zettelkasten Self-Organizing Memory}

Inspired by the A-MEM framework, the Zettelkasten module implements
a self-organizing note graph where each \texttt{MemoryNote} carries:

\begin{itemize}[leftmargin=*]
  \item Auto-extracted keywords (stop-word filtered, max 20)
  \item Embedding vector for semantic similarity
  \item Bidirectional links to related notes
  \item Memory strength $S$ (Ebbinghaus)
  \item Access count and temporal metadata
\end{itemize}

\paragraph{The A-MEM Insight: Keyword Evolution.}
When a new note $n_\text{new}$ is added, the system:
(1)~finds semantically similar notes by cosine similarity above
threshold $\tau_\text{link}$ (default 0.70);
(2)~creates bidirectional links;
(3)~\emph{evolves existing notes' keywords} by injecting up to 3
novel keywords from $n_\text{new}$. This causes the knowledge graph
to self-organize---older notes become richer as related knowledge
accumulates.

\paragraph{Ebbinghaus Forgetting Curve.}
Each note's retention probability follows:

\begin{equation}
R = e^{-t / (S \cdot 24)}
\label{eq:ebbinghaus}
\end{equation}

where $t$ is hours since last access and $S$ is memory strength
(starts at 1.0, increases by 0.5 per access, capped at 10.0).
Notes with $R < \tau_\text{forget}$ (default 0.2) are candidates
for removal. The forgetting pass also cleans up indexes and
removes dangling links.

\subsection{RL-Guided Memory Evolution}

The \texttt{MemoryEvolution} engine runs periodic cycles with
five steps:

\begin{enumerate}[leftmargin=*]
  \item \textbf{Composite Scoring}: Each item receives an importance
    score combining four signals:
    \begin{equation}
    I = 0.3 \cdot f_\text{freq} + 0.3 \cdot f_\text{recency}
      + 0.25 \cdot f_\text{conn} + 0.15 \cdot f_\text{salience}
    \end{equation}
    where $f_\text{freq} = \ln(\text{accesses}+1)/10$,
    $f_\text{recency} = 1/(1+h/168)$ (1-week half-life),
    $f_\text{conn} = \min(\text{connectivity}, 20)/20$,
    and $f_\text{salience} = \min(\text{words}, 50)/50$.

  \item \textbf{Temporal Decay}: Importance is multiplied by a
    decay factor $D$ according to the configured function:
    exponential ($e^{-\lambda t}$), linear ($1 - \lambda t$),
    or step (0 if $t > \tau$).

  \item \textbf{RL Keep/Forget}: Items below the forget threshold
    (default 0.1) are evaluated by the Q-table RL agent. State is
    discretized into bins of (importance, access frequency, recency,
    connectivity). The agent uses $\epsilon$-greedy exploration
    ($\epsilon=0.1$) with Q-learning updates:
    \begin{equation}
    Q(s,a) \leftarrow Q(s,a) + \alpha[r + \gamma \max_{a'} Q(s',a')
      - Q(s,a)]
    \end{equation}
    with $\alpha=0.1$, $\gamma=0.95$.

  \item \textbf{Consolidation}: Items with cosine similarity $> 0.95$
    are merged, keeping the higher-importance item.

  \item \textbf{Overflow Eviction}: If items exceed
    \texttt{max\_items} (default 100K), the lowest-importance items
    are truncated.
\end{enumerate}

% ════════════════════════════════════════════════════════════════════════════
\section{Multi-Graph Memory (MAGMA)}
\label{sec:multigraph}

\subsection{Four Orthogonal Graph Views}

Inspired by the MAGMA architecture, each \texttt{MemoryItem} is
indexed across four independent graph views:

\begin{enumerate}[leftmargin=*]
  \item \textbf{Semantic Graph}: Edges represent embedding cosine
    similarity above a threshold (default 0.75). Uses brute-force
    $O(n^2)$ comparison with planned HNSW upgrade.
  \item \textbf{Temporal Graph}: Edges encode Allen's interval algebra
    relations (before, after, during, overlaps, etc.) between items'
    validity intervals.
  \item \textbf{Causal Graph}: Directed edges represent cause-effect
    relationships extracted via the \texttt{CausalExtractor}, which
    applies five regex patterns with typed confidence scores:
    \emph{caused/led to} (0.8), \emph{because/since} (0.7),
    \emph{therefore/consequently} (0.75), \emph{if-then} (0.6),
    and \emph{enabled/allowed} (0.7). Edges below the configurable
    threshold (default 0.6) are discarded. The graph supports both
    forward chaining (\texttt{traverse()}: BFS with per-hop
    confidence decay) and backward chaining (\texttt{root\_causes()}:
    enumerates all root-cause paths up to a configurable depth).
  \item \textbf{Entity Graph}: Standard typed entity relationships
    extracted from content.
\end{enumerate}

Ingestion is parallelized via \texttt{tokio::join!}, concurrently
indexing each item across all four graphs.

\subsection{Intent-Aware Query Router}

The \texttt{IntentRouter} classifies queries into eight intent types
using keyword-scoring heuristics:

\begin{table}[H]
\centering
\footnotesize
\caption{Intent types and primary graph views.}
\label{tab:intents}
\begin{tabularx}{\columnwidth}{l@{\hspace{4pt}}X@{\hspace{4pt}}l}
\toprule
\textbf{Intent} & \textbf{Example} & \textbf{View} \\
\midrule
Factual & Who is Alice? & Entity \\
Temporal & What happened in 2024? & Temporal \\
Causal & Why did X fail? & Causal \\
Hypothetical & What if X hadn't\ldots? & Causal \\
Comparative & Compare X and Y & Entity \\
Aggregative & How many at X? & Entity \\
MultiHop & Alice's boss reports to? & Entity \\
Exploratory & Tell me about X & Semantic \\
\bottomrule
\end{tabularx}
\end{table}

Each intent maps to a \texttt{TraversalPlan} specifying view weights,
max hops (2--5), max nodes per graph (20--40), and fusion strategy
(RRF with $k=60$ or weighted sum).

\subsection{Fusion Engine}

The \texttt{FusionEngine} merges \texttt{SubgraphResult}s from
multiple views using type-aligned RRF with cross-view reinforcement.
Items appearing in multiple views receive additive score boosts
proportional to the view weights specified by the intent router.

% ════════════════════════════════════════════════════════════════════════════
\section{Temporal Subsystem}
\label{sec:temporal}

\subsection{Bi-Temporal Model}

Every \texttt{Relationship} in NeuroGraph carries two temporal
intervals:
\begin{itemize}[leftmargin=*]
  \item \texttt{valid\_from} / \texttt{valid\_until}: The real-world
    validity period of the fact.
  \item \texttt{created\_at}: When the fact was ingested (system time).
\end{itemize}

The \texttt{TemporalManager} provides three temporal operations:
\begin{itemize}[leftmargin=*]
  \item \texttt{snapshot\_at(t)}: Returns all entities and relationships
    valid at timestamp $t$, filtering by \texttt{valid\_from $\leq t$}
    and \texttt{valid\_until} is null or $> t$.
  \item \texttt{what\_changed(t1, t2)}: Returns entities added,
    modified, and relationships invalidated between $t_1$ and $t_2$.
  \item \texttt{build\_timeline()}: Constructs a chronological event
    list for visualization (G6 Timebar integration).
\end{itemize}

\subsection{Git-Like Branch Isolation}

The \texttt{BranchManager} implements a copy-on-write (COW) branching
model for multi-agent memory isolation:

\begin{itemize}[leftmargin=*]
  \item \textbf{Create}: Lightweight branch forked from a parent,
    storing only a delta set of added/removed entity and relationship
    IDs. Thread-safe via \texttt{DashMap}.
  \item \textbf{Read}: Visibility check walks the branch chain
    (branch delta $\to$ parent delta $\to$ root), implementing
    COW semantics.
  \item \textbf{Diff}: Computes the symmetric difference between
    two branches' delta sets.
  \item \textbf{Merge}: Applies source branch's deltas to target,
    detecting conflicts (e.g., entity added in source but removed
    in target). Conflicts are flagged but the merge proceeds.
  \item \textbf{Abandon}: Marks a branch as closed without merging.
\end{itemize}

Configuration enforces a maximum of 100 open branches and optional
auto-abandonment after 24 hours.

\subsection{Logical Clock and Versioning}

The \texttt{LogicalClock} implements a hybrid logical clock combining
wall time (\texttt{DateTime<Utc>}) with a process-global monotonic
counter (\texttt{AtomicU64}). The \texttt{happened\_before()} method
compares wall times first and breaks ties with the counter, ensuring
deterministic ordering even under concurrent ingestion.

The versioning module provides \texttt{FactVersionChain}---a
chronological list of \texttt{FactVersion} records for a single
relationship. Each version carries the fact text, validity interval,
confidence, and an optional \texttt{superseded\_by} link. The
\texttt{version\_at(t)} method returns the fact that was valid at
any historical timestamp.

\texttt{EntityHistory} tracks modification records per entity,
supporting six change types: \texttt{Created}, \texttt{SummaryUpdated},
\texttt{Merged} (deduplication), \texttt{TypeChanged},
\texttt{AttributesUpdated}, and \texttt{DecayMarked}. The
\texttt{changes\_between(t1, t2)} method enables audit trail queries.

\subsection{Graph-Level Forgetting Engine}

The \texttt{temporal::forgetting} module provides a separate,
graph-level forgetting mechanism distinct from the memory-layer
evolution engine (\S\ref{sec:memory}). The \texttt{ForgettingEngine}
operates directly on \texttt{Entity} nodes in the \texttt{GraphDriver}:

\begin{equation}
I_\text{decay}(e) = I_0(e) \cdot e^{-\lambda d}
  + \frac{\ln(a+1)}{10}
  + \frac{\ln(r+1)}{10}
\end{equation}

where $I_0$ is the entity's base importance, $\lambda$ is the
\texttt{daily\_decay\_rate} (default 0.01), $d$ is days since last
access, $a$ is the access count (log-scaled boost), and $r$ is
the relationship count (connectivity boost). Entities with
$I_\text{decay} < \tau_\text{prune}$ (default 0.1) and access
count $<$ \texttt{min\_access\_count} (default 5) become prune
candidates.

The \texttt{decay\_pass(auto\_prune, group\_id)} method evaluates
all entities, updates importance scores, and optionally deletes
candidates in batches (\texttt{max\_prune\_batch}, default 100).
A separate TTL-based expiration via \texttt{is\_expired()} removes
entities exceeding \texttt{default\_ttl} regardless of importance.

% ════════════════════════════════════════════════════════════════════════════
\section{Community Detection}
\label{sec:community}

NeuroGraph implements two community detection algorithms:

\begin{itemize}[leftmargin=*]
  \item \textbf{Louvain}: Greedy modularity optimization with
    configurable resolution parameter, max iterations, and minimum
    community size.
  \item \textbf{Leiden}: Refinement phase on top of Louvain for
    higher-quality, well-connected communities.
\end{itemize}

Both operate on the entity-relationship graph extracted from the
\texttt{GraphDriver}. Detected communities are stored as
\texttt{Community} objects with member entity IDs, hierarchical
level, and optional LLM-generated summaries.

The \texttt{IncrementalCommunityUpdater} supports post-ingestion
updates: given a set of newly added entity IDs, it re-evaluates
only the affected neighborhoods rather than rerunning detection
on the full graph.

The \texttt{CommunitySummarizer} generates natural-language summaries
via LLM (with rule-based fallback), enabling the DRIFT global
search path to use community summaries as map-reduce intermediaries.

% ════════════════════════════════════════════════════════════════════════════
\section{Research Paper Intelligence}
\label{sec:research}

\subsection{PDF Ingestion Pipeline}

The PDF module provides a two-tier parsing architecture:

\begin{itemize}[leftmargin=*]
  \item \textbf{Fast mode}: Uses \texttt{pdf-extract} for rapid
    text extraction (~1--5ms per page). Suitable for bulk ingestion.
  \item \textbf{Structured mode}: Heuristic-based academic section
    detection that identifies title, authors, abstract, sections
    (Introduction, Methods, Results, etc.), references, equations,
    and figures via regex patterns on formatting cues.
  \item \textbf{Auto mode}: The \texttt{DocumentClassifier} analyzes
    the raw text to estimate column count, detect academic indicators
    (abstract, references, affiliations), and check for mathematical
    symbols, then recommends the optimal strategy.
\end{itemize}

The \texttt{SmartChunker} splits text into semantically coherent
chunks with configurable maximum tokens (default 512) and sentence
overlap (default 2), respecting section boundaries.

\subsection{Multi-Source Academic Search}

The \texttt{papers} module provides unified search across three
academic APIs through the \texttt{UnifiedPaperSearch} aggregator:

\begin{table}[H]
\centering
\footnotesize
\caption{Academic search API integrations.}
\label{tab:papers}
\begin{tabularx}{\columnwidth}{lXX}
\toprule
\textbf{Source} & \textbf{API} & \textbf{Features} \\
\midrule
arXiv & Atom XML feed & Categories, PDFs \\
Semantic Scholar & REST v1 & Citations, refs \\
PubMed & E-utilities & MeSH, PMIDs \\
\bottomrule
\end{tabularx}
\end{table}

Results are deduplicated by title similarity and can be sorted
by relevance, citation count, or date. The CLI supports
\texttt{--since} year filtering and \texttt{--download} for
automatic PDF retrieval.

\subsection{RAG Chat Engine}

The chat module implements a complete conversational RAG pipeline:

\begin{itemize}[leftmargin=*]
  \item \textbf{Context Assembly}: The \texttt{ContextBuilder}
    retrieves the top-$k$ relevant chunks from the knowledge
    graph, formats them with source citations (paper title,
    section, page number).
  \item \textbf{Conversation History}: The \texttt{ConversationHistory}
    maintains a sliding window of messages with role annotations
    (User, Assistant, System) and attached source citations.
  \item \textbf{RAG Pipeline}: Formats a system prompt incorporating
    retrieved context, conversation history, and applicable
    procedural rules from L4 memory.
  \item \textbf{REPL}: Terminal-based interactive loop with
    \texttt{/clear}, \texttt{/quit} commands and colored output.
\end{itemize}

% ════════════════════════════════════════════════════════════════════════════
\section{REST API and Dashboard}
\label{sec:server}

\subsection{API Design}

The server module provides a RESTful API built on Axum 0.8 with
Tower middleware for CORS and tracing:

\begin{table}[H]
\centering
\scriptsize
\caption{REST API endpoints.}
\label{tab:api}
\begin{tabular}{@{}llp{2.2cm}@{}}
\toprule
\textbf{Verb} & \textbf{Endpoint} & \textbf{Description} \\
\midrule
\textsf{GET}  & \texttt{/health} & Health check \\
\textsf{GET}  & \texttt{/stats} & Graph statistics \\
\textsf{POST} & \texttt{/papers/ingest} & PDF upload \\
\textsf{POST} & \texttt{/papers/ingest-text} & Text ingestion \\
\textsf{POST} & \texttt{/papers/ingest-url} & URL ingestion \\
\textsf{GET}  & \texttt{/papers} & List papers \\
\textsf{POST} & \texttt{/query} & NL query \\
\textsf{GET}  & \texttt{/graph} & Graph overview \\
\textsf{POST} & \texttt{/chat} & RAG chat \\
\textsf{POST} & \texttt{/search/papers} & Paper search \\
\textsf{POST} & \texttt{/search/entities} & Entity search \\
\bottomrule
\end{tabular}
\vspace{2pt}

\raggedright\scriptsize All paths prefixed with \texttt{/api/v1}.
\end{table}

All responses follow a uniform \texttt{ApiResponse\{success, data,
error\}} envelope. The server state (\texttt{AppState}) wraps a
shared \texttt{NeuroGraph} instance, paper metadata, and
conversation histories using \texttt{parking\_lot::RwLock} for
synchronization.

\subsection{Embedded Dashboard}

Static dashboard assets are embedded in the binary via
\texttt{rust-embed}, enabling a single-binary deployment with no
external file dependencies. The dashboard provides a G6-based
graph visualization with node type encoding, importance sizing,
and memory tier coloring.

\subsection{Model Context Protocol (MCP)}
\label{sec:mcp}

The \texttt{neurograph-mcp} crate exposes NeuroGraph as an MCP
server, enabling integration with Claude, Cursor, and VS~Code
Copilot. The MCP server wraps the \texttt{neurograph-core} API
and provides tool definitions for \texttt{add\_text},
\texttt{query}, \texttt{search\_papers}, \texttt{entity\_history},
and \texttt{stats}. Transport is supported over both \texttt{stdio}
(for IDE integration) and Server-Sent Events (SSE) for web-based
clients.

% ════════════════════════════════════════════════════════════════════════════
\section{Embedding Subsystem}
\label{sec:embedders}

NeuroGraph supports four embedding providers behind the
\texttt{Embedder} trait:

\begin{enumerate}[leftmargin=*]
  \item \textbf{HashEmbedder} (default): Deterministic 128-dim
    vectors via SeaHash feature hashing. Zero-cost, no API keys,
    suitable for structural deduplication.
  \item \textbf{OpenAI Embedder}: \texttt{text-embedding-3-small}
    via the \texttt{async-openai} crate.
  \item \textbf{Ollama Embedder}: Local models (e.g.,
    \texttt{nomic-embed-text}) via Ollama's HTTP API.
  \item \textbf{Cached Embedder}: LRU cache wrapper (default
    capacity 10,000) around any provider, keyed by BLAKE3 content
    hash.
\end{enumerate}

The \texttt{EmbeddingRouter} selects the appropriate provider
based on configuration and falls back gracefully if a provider
is unavailable.

% ════════════════════════════════════════════════════════════════════════════
\section{Command-Line Interface}
\label{sec:cli}

The CLI (\texttt{neurograph-cli}) provides 12 subcommands:

\begin{lstlisting}[language=bash,style={}]
neurograph ingest --pdf paper.pdf
neurograph ingest --dir ./papers/ --strategy auto
neurograph ingest --url https://arxiv.org/pdf/...
neurograph query "Who works at Anthropic?" --at 2024
neurograph search "attention" --arxiv --s2 --since 2023
neurograph chat --model ollama:llama3.2
neurograph dashboard --port 8000 --open
neurograph serve --port 8000
neurograph stats
neurograph history "Alice"
neurograph what-changed 2024-01-01 2025-01-01
neurograph communities
neurograph bench
neurograph info
\end{lstlisting}

The CLI uses \texttt{clap} 4 for argument parsing, \texttt{colored}
for terminal output, and \texttt{indicatif} for progress bars during
bulk ingestion.

% ════════════════════════════════════════════════════════════════════════════
\section{Deployment and Infrastructure}
\label{sec:deployment}

\subsection{Build and Distribution}

NeuroGraph supports multiple distribution channels:
\begin{itemize}[leftmargin=*]
  \item \textbf{Cargo}: \texttt{cargo install neurograph-cli}
  \item \textbf{Docker}: Multi-stage Dockerfile with Rust builder
    and minimal runtime image
  \item \textbf{Docker Compose}: Full stack with Neo4j, Ollama,
    and dashboard
\end{itemize}

\subsection{Security and Compliance}

The project implements comprehensive security practices:
\texttt{cargo-deny} for dependency auditing, OSSF Security Insights
YAML, a responsible AI transparency document, and a detailed
SECURITY.md with vulnerability disclosure procedures.

% ════════════════════════════════════════════════════════════════════════════
\section{Related Work}
\label{sec:related}

NeuroGraph draws inspiration from and extends several systems:

\begin{itemize}[leftmargin=*]
  \item \textbf{Microsoft GraphRAG} \cite{graphrag2024}: Hierarchical
    Leiden community detection and global/local search. NeuroGraph
    extends this with DRIFT adaptive switching and five-strategy
    hybrid retrieval.
  \item \textbf{Graphiti} \cite{graphiti2024}: Temporal knowledge
    graph with episode-based ingestion. NeuroGraph adds bi-temporal
    validity, branch isolation, and RL-guided forgetting.
  \item \textbf{Cognee}: Pipeline-oriented knowledge processing.
    NeuroGraph adopts the clean \texttt{add/query} API pattern.
  \item \textbf{LightRAG}: Lightweight RAG with knowledge graphs.
    NeuroGraph provides richer retrieval and memory management.
  \item \textbf{Mem0/EverMemOS}: Tiered memory with L1--L4
    hierarchy. NeuroGraph implements the tiered model and adds
    RL-guided evolution.
  \item \textbf{A-MEM}: Zettelkasten self-organizing memory.
    NeuroGraph implements the keyword evolution insight with
    Ebbinghaus forgetting integration.
  \item \textbf{MAGMA}: Multi-graph memory architecture.
    NeuroGraph implements four orthogonal views with intent-aware
    routing and RRF fusion.
\end{itemize}

% ════════════════════════════════════════════════════════════════════════════
\section{Conclusion and Future Work}
\label{sec:conclusion}

NeuroGraph demonstrates that a single, cohesive system can integrate
temporal knowledge management, advanced hybrid retrieval,
biologically-inspired memory evolution, and multi-agent collaboration.
By implementing the entire stack in Rust, we achieve both developer
ergonomics (via the builder pattern and zero-config defaults) and
production-grade performance (async I/O, lock-free concurrency).

\paragraph{Future Directions.}
\begin{itemize}[leftmargin=*]
  \item \textbf{HNSW Vector Index}: Replace brute-force similarity
    search with approximate nearest neighbors for sub-linear scaling.
  \item \textbf{Neo4j Driver}: Full Cypher-backed persistent storage
    for enterprise deployments.
  \item \textbf{WASM/Python Bindings}: Cross-language SDKs via FFI
    and WebAssembly compilation.
  \item \textbf{Agentic Framework}: Multi-agent orchestration with
    role-based branch policies and supervisor approval workflows.
  \item \textbf{Continuous Learning}: Online RL policy updates from
    user retrieval feedback, closing the forgetting optimization loop.
\end{itemize}

% ════════════════════════════════════════════════════════════════════════════
\begin{thebibliography}{9}

\bibitem{graphrag2024}
Edge, D. \textit{et al.}
\newblock From Local to Global: A Graph RAG Approach to Query-Focused
  Summarization.
\newblock \textit{arXiv preprint arXiv:2404.16130}, 2024.

\bibitem{graphiti2024}
Zep AI.
\newblock Graphiti: Building Real-Time, Temporally Aware Knowledge Graphs.
\newblock \textit{GitHub: getzep/graphiti}, 2024.

\bibitem{amem2025}
Xu, Z. \textit{et al.}
\newblock A-MEM: Agentic Memory for LLM Agents.
\newblock \textit{arXiv preprint}, 2025.

\bibitem{evermemos2026}
Zhang, L. \textit{et al.}
\newblock EverMemOS: Tiered Memory Operating System for Lifelong AI Agents.
\newblock \textit{arXiv preprint}, 2026.

\bibitem{drift2026}
Microsoft Research.
\newblock DRIFT Search: Dynamic Reasoning and Inference with Flexible Traversal.
\newblock \textit{GraphRAG v3 Documentation}, 2026.

\bibitem{lightrag2024}
Guo, Z. \textit{et al.}
\newblock LightRAG: Simple and Fast Retrieval-Augmented Generation.
\newblock \textit{arXiv preprint}, 2024.

\bibitem{rrf2009}
Cormack, G.V., Clarke, C.L.A., and Buettcher, S.
\newblock Reciprocal Rank Fusion Outperforms Condorcet and Individual Rank
  Learning Methods.
\newblock \textit{Proc. SIGIR}, 2009.

\bibitem{leiden2019}
Traag, V.A., Waltman, L., and van Eck, N.J.
\newblock From Louvain to Leiden: guaranteeing well-connected communities.
\newblock \textit{Scientific Reports}, 9(1):5233, 2019.

\bibitem{ebbinghaus1885}
Ebbinghaus, H.
\newblock \"{U}ber das Ged\"{a}chtnis.
\newblock Duncker \& Humblot, Leipzig, 1885.

\end{thebibliography}

\end{document}
